\documentclass[10pt,a4paper]{amsart}
\usepackage{amsmath,amssymb,amsthm,graphicx}
% Retained minimal preamble for the Stacks Project chapter more-algebra.tex (More on Algebra).
% The chapter file pinned-original-source.tex carries only \input{preamble} on its first line, so no
% preamble text lives inside the chapter itself.  The Stacks Project uses one single shared preamble for
% every chapter; the snapshot stages it identically as sources/stacks-chow/preamble.tex and
% sources/stacks-etale-cohomology/preamble.tex.  The block below is copied mechanically, line for line,
% from sources/stacks-chow/preamble.tex lines 14-18 (the Xy-pic setup) and 154-217 (the native theorem
% declarations together with the mathematical macros).  Nothing was added, removed or redefined; the
% wrapper class and the amsmath/amssymb/amsthm/graphicx base packages are supplied by the assembler.
% Chapter: more-algebra.tex -- More on Algebra.
% For commutative diagrams we use Xy-pic
\usepackage[all]{xy}

% We use 2cell for 2-commutative diagrams.
\xyoption{2cell}

% Theorem environments.
%
\theoremstyle{plain}
\newtheorem{theorem}[subsection]{Theorem}
\newtheorem{proposition}[subsection]{Proposition}
\newtheorem{lemma}[subsection]{Lemma}

\theoremstyle{definition}
\newtheorem{definition}[subsection]{Definition}
\newtheorem{example}[subsection]{Example}
\newtheorem{exercise}[subsection]{Exercise}
\newtheorem{situation}[subsection]{Situation}

\theoremstyle{remark}
\newtheorem{remark}[subsection]{Remark}
\newtheorem{remarks}[subsection]{Remarks}

\numberwithin{equation}{subsection}

% Macros
%
\def\lim{\mathop{\mathrm{lim}}\nolimits}
\def\colim{\mathop{\mathrm{colim}}\nolimits}
\def\Spec{\mathop{\mathrm{Spec}}}
\def\Hom{\mathop{\mathrm{Hom}}\nolimits}
\def\Ext{\mathop{\mathrm{Ext}}\nolimits}
\def\SheafHom{\mathop{\mathcal{H}\!\mathit{om}}\nolimits}
\def\SheafExt{\mathop{\mathcal{E}\!\mathit{xt}}\nolimits}
\def\Sch{\mathit{Sch}}
\def\Mor{\mathop{\mathrm{Mor}}\nolimits}
\def\Ob{\mathop{\mathrm{Ob}}\nolimits}
\def\Sh{\mathop{\mathit{Sh}}\nolimits}
\def\NL{\mathop{N\!L}\nolimits}
\def\CH{\mathop{\mathrm{CH}}\nolimits}
\def\proetale{{pro\text{-}\acute{e}tale}}
\def\etale{{\acute{e}tale}}
\def\QCoh{\mathit{QCoh}}
\def\Ker{\mathop{\mathrm{Ker}}}
\def\Im{\mathop{\mathrm{Im}}}
\def\Coker{\mathop{\mathrm{Coker}}}
\def\Coim{\mathop{\mathrm{Coim}}}

% Boxtimes
%
\DeclareMathSymbol{\boxtimes}{\mathbin}{AMSa}{"02}

%
% Macros for moduli stacks/spaces
%
\def\QCohstack{\mathcal{QC}\!\mathit{oh}}
\def\Cohstack{\mathcal{C}\!\mathit{oh}}
\def\Spacesstack{\mathcal{S}\!\mathit{paces}}
\def\Quotfunctor{\mathrm{Quot}}
\def\Hilbfunctor{\mathrm{Hilb}}
\def\Curvesstack{\mathcal{C}\!\mathit{urves}}
\def\Polarizedstack{\mathcal{P}\!\mathit{olarized}}
\def\Complexesstack{\mathcal{C}\!\mathit{omplexes}}
% \Pic is the operator that assigns to X its picard group, usage \Pic(X)
% \Picardstack_{X/B} denotes the Picard stack of X over B
% \Picardfunctor_{X/B} denotes the Picard functor of X over B
\def\Pic{\mathop{\mathrm{Pic}}\nolimits}
\def\Picardstack{\mathcal{P}\!\mathit{ic}}
\def\Picardfunctor{\mathrm{Pic}}
\def\Deformationcategory{\mathcal{D}\!\mathit{ef}}

% --- appended 2026-10-12 (Chief): the project own definitions of the reference/slogan/history
% display environments, copied mechanically from the shared Stacks preamble lines 8-12
% (\usepackage{verbatim} plus the three \newenvironment lines, which the chapter preamble keeps
% verbatim from the project preamble). Nothing added, removed or redefined: these environments
% carry no mathematical content and the project itself typesets them as comments. Without these
% three lines a counted statement whose own text uses them fails with Environment slogan undefined.
\usepackage{verbatim}
\newenvironment{reference}{\comment}{\endcomment}
\newenvironment{slogan}{\comment}{\endcomment}
\newenvironment{history}{\comment}{\endcomment}
\usepackage{cleveref}
\crefname{originallemma}{original lemma}{original lemmas}
\Crefname{originallemma}{Original lemma}{Original lemmas}
\makeatletter
\newlabel{lemma-flat-finite-presentation-perfect}{{067J}{source}{}{sourcepointer-1}{}}
\newlabel{lemma-flat-finite-presentation-perfect@cref}{{[originallemma][0][]067J}{[1][0][]source}{}{}{}}
\makeatother
\crefname{originallemma}{original lemma}{original lemmas}
\Crefname{originallemma}{Original lemma}{Original lemmas}
\makeatletter
\newlabel{lemma-relative-pseudo-coherent-is-moot}{{0DHQ}{source}{}{sourcepointer-2}{}}
\newlabel{lemma-relative-pseudo-coherent-is-moot@cref}{{[originallemma][0][]0DHQ}{[1][0][]source}{}{}{}}
\makeatother
\crefname{originallemma}{original lemma}{original lemmas}
\Crefname{originallemma}{Original lemma}{Original lemmas}
\makeatletter
\newlabel{lemma-composition-relative-pseudo-coherent}{{067D}{source}{}{sourcepointer-3}{}}
\newlabel{lemma-composition-relative-pseudo-coherent@cref}{{[originallemma][0][]067D}{[1][0][]source}{}{}{}}
\makeatother
\crefname{originaldefinition}{original definition}{original definitions}
\Crefname{originaldefinition}{Original definition}{Original definitions}
\makeatletter
\newlabel{definition-relatively-pseudo-coherent}{{065I}{source}{}{sourcepointer-4}{}}
\newlabel{definition-relatively-pseudo-coherent@cref}{{[originaldefinition][0][]065I}{[1][0][]source}{}{}{}}
\makeatother
\title{Coincidence of absolute and relative pseudo-coherence for flat finite presentation quotients}
\author{The Stacks Project Authors}
\date{stacks-project:more-algebra.tex:local-snapshot-20261012; GNU FDL 1.2 or later}
\begin{document}
\maketitle
\setcounter{subsection}{1}
\noindent\textit{Editorial note.} All delivered mathematical text is the author's own, copied from the pinned Stacks Project chapter file more-algebra.tex; no statement or proof step was rewritten, reordered or removed. The counted span is the author environment and its own proof environment. References inside the span are handed to the assembler as navigation stubs that carry the canonical tag of the cited tag only; whether each of them is load-bearing is recorded on the stub itself.\par\medskip
\noindent\textit{Allowed standard prerequisites (editorial).}\begin{itemize}
\item Flat and finite presentation ring maps are perfect, hence pseudo-coherent ring maps. invoked at proof lines 21901-21903. Makes relative pseudo-coherence over R moot.
\item Over pseudo-coherent ring maps absolute and relative pseudo-coherence agree. invoked at proof lines 21904-21905. Gives (1) iff (2) and (4) iff (5).
\item Relative pseudo-coherence is transitive under composition, and A is pseudo-coherent relative to R. invoked at proof lines 21908-21910. Gives (2) iff (3).
\item The definition of relative pseudo-coherence over a surjective ring map reduces to absolute pseudo-coherence of the restriction. used at proof lines 21911-21913. Gives (3) iff (4).
\end{itemize}\medskip
\ifdefined\hypertarget\hypertarget{sourcepointer-1}{}\fi
\noindent\textit{Source reference.} Original lemma 067J: sources/stacks-more-algebra/pinned-original-source.tex, source lines 21815--21815. Cited at proof lines 21901-21902 to identify the ring maps as perfect; the equivalence assembly is inside the counted span.\par\medskip
\ifdefined\hypertarget\hypertarget{sourcepointer-2}{}\fi
\noindent\textit{Source reference.} Original lemma 0DHQ: sources/stacks-more-algebra/pinned-original-source.tex, source lines 21870--21870. Cited at proof lines 21904-21905 for (1) iff (2) and (4) iff (5); the remaining equivalences are inside the counted span.\par\medskip
\ifdefined\hypertarget\hypertarget{sourcepointer-3}{}\fi
\noindent\textit{Source reference.} Original lemma 067D: sources/stacks-more-algebra/pinned-original-source.tex, source lines 21638--21638. Cited at proof lines 21908-21910 for (2) iff (3); the surjective-case equivalence is inside the counted span.\par\medskip
\ifdefined\hypertarget\hypertarget{sourcepointer-4}{}\fi
\noindent\textit{Source reference.} Original definition 065I: sources/stacks-more-algebra/pinned-original-source.tex, source lines 21332--21332. Cited at proof lines 21912-21912 for the direct equivalence (3) iff (4) under the surjection; the assembly is inside the counted span.\par\medskip
\begin{lemma}
\label{lemma-more-relative-pseudo-coherent-is-moot}
Let $R \to B \to A$ be ring maps with $\varphi : B \to A$ surjective and
$R \to B$ and $R \to A$ flat and of finite presentation. For $K \in D(A)$
denote $\varphi_*K \in D(B)$ the restriction.
The following are equivalent
\begin{enumerate}
\item $K$ is pseudo-coherent,
\item $K$ is pseudo-coherent relative to $R$,
\item $K$ is pseudo-coherent relative to $A$,
\item $\varphi_*K$ is pseudo-coherent,
\item $\varphi_*K$ is pseudo-coherent relative to $R$.
\end{enumerate}
Similar holds for $m$-pseudo-coherence.
\end{lemma}

\begin{proof}
Observe that $R \to A$ and $R \to B$ are perfect ring maps
(Lemma \ref{lemma-flat-finite-presentation-perfect})
hence a fortiori pseudo-coherent ring maps.
Thus (1) $\Leftrightarrow$ (2) and (4) $\Leftrightarrow$ (5)
by Lemma \ref{lemma-relative-pseudo-coherent-is-moot}.

\medskip\noindent
Using that $A$ is pseudo-coherent relative to $R$ we use
Lemma \ref{lemma-composition-relative-pseudo-coherent}
to see that (2) $\Leftrightarrow$ (3).
However, since $A \to B$ is surjective, we see directly from
Definition \ref{definition-relatively-pseudo-coherent}
that (3) is equivalent with (4).
\end{proof}

\end{document}
